\documentclass[11pt, a4paper]{article}
\usepackage[margin=2.5cm]{geometry}
\usepackage{amsmath, amssymb}
\usepackage{booktabs}
\usepackage{array}

\begin{document}

\subsection{Extension: Population Balance Analysis and Improvement (C2)}

\subsubsection{Population Assignment and Baseline Evaluation}

The spectral partition obtained from the unweighted Fiedler vector in Figure~1.22 relies solely on shared boundary adjacency and completely disregards voter population. Consequently, it cannot guarantee the legal redistricting requirement of \textit{equity} (equal population across districts). To analyze this limitation, we assign synthetic yet realistic residential populations to the 16 census blocks (2000--2015), yielding a total population of $P_{\text{total}} = 2000$ residents and an ideal population per district of $P_{\text{ideal}} = P_{\text{total}} / 2 = 1000$.

Table~\ref{tab:c2_block_pop} lists the assigned population $p_i$ for each census block alongside its baseline district assignment from Figure~1.22.\footnote{Erratum: On page~35 of Arangala (2023), the text erroneously includes census block~2016 in $G_1$. Because the study graph only encompasses blocks~2000 through~2015 (16 vertices in total), block~2016 is excluded in our corrected partition.}

\begin{table}[htbp]
\centering
\caption{Assigned population for census blocks 2000--2015 and baseline Fiedler partition from Figure~1.22.}
\label{tab:c2_block_pop}
\begin{tabular}{@{}cccc|cccc@{}}
\toprule
Block & District & Boundary & Pop. ($p_i$) & Block & District & Boundary & Pop. ($p_i$) \\
\midrule
2000 & $G_1$ & No  & 130 & 2008 & $G_2$ & No  & 160 \\
2001 & $G_1$ & No  & 80  & 2009 & $G_2$ & No  & 130 \\
2002 & $G_1$ & No  & 75  & 2010 & $G_1$ & Yes & 180 \\
2003 & $G_1$ & No  & 115 & 2011 & $G_1$ & No  & 110 \\
2004 & $G_1$ & Yes & 260 & 2012 & $G_1$ & No  & 55  \\
2005 & $G_2$ & Yes & 155 & 2013 & $G_1$ & No  & 125 \\
2006 & $G_2$ & Yes & 135 & 2014 & $G_1$ & No  & 80  \\
2007 & $G_2$ & No  & 140 & 2015 & $G_1$ & No  & 70  \\
\midrule
\multicolumn{4}{l|}{\textbf{Baseline $G_1$ (11 blocks):} $P(G_1) = 1280$} & \multicolumn{4}{l}{\textbf{Baseline $G_2$ (5 blocks):} $P(G_2) = 720$} \\
\bottomrule
\end{tabular}
\end{table}

To quantify district equity rigorously, we define two standard population balance measures: the \textit{maximum relative deviation} $\delta_{\max}$ from the ideal district population (Equation~\ref{eq:delta_max}) and the \textit{overall population deviation range} $\Delta$ (Equation~\ref{eq:delta_range}):
\begin{equation}
\delta_{\max} = \max_{k \in \{1, 2\}} \frac{|P(G_k) - P_{\text{ideal}}|}{P_{\text{ideal}}} \times 100\%,
\label{eq:delta_max}
\end{equation}
\begin{equation}
\Delta = \frac{\max_k P(G_k) - \min_k P(G_k)}{P_{\text{ideal}}} \times 100\%.
\label{eq:delta_range}
\end{equation}
In U.S. redistricting practice, an overall deviation range of $\Delta \le 10\%$ is widely adopted as a standard equity threshold for state and local districts.

For the baseline Fiedler partition, summing the populations in Table~\ref{tab:c2_block_pop} yields $P(G_1) = 1280$ and $P(G_2) = 720$. Substituting these values into Equations~\eqref{eq:delta_max} and~\eqref{eq:delta_range} gives:
\[
\delta_{\max} = \frac{|1280 - 1000|}{1000} \times 100\% = 28.000\%, \qquad \Delta = \frac{1280 - 720}{1000} \times 100\% = 56.000\%.
\]
\textit{Interpretation:} Because $G_1$ contains 11 blocks while $G_2$ contains only 5, the baseline Fiedler cut exhibits a severe population imbalance of 560 residents ($\Delta = 56.000\%$), far exceeding the $10\%$ benchmark despite achieving a minimal cut size of only 3 boundary edges: $(2004, 2005)$, $(2004, 2006)$, and $(2010, 2005)$.

\subsubsection{Improving Balance via Boundary-Block Reassignment}

To improve equity without breaking district \textit{contiguity}, we apply a greedy single boundary-block reassignment procedure:
\begin{enumerate}
\item Identify all boundary vertices in the overpopulated district $G_1$ that share at least one edge with $G_2$. From the adjacency matrix of Figure~1.22, the only candidate boundary vertices are block~2004 (adjacent to 2005 and 2006) and block~2010 (adjacent to 2005).
\item For each candidate $u \in \{2004, 2010\}$, check whether both $G_1 \setminus \{u\}$ and $G_2 \cup \{u\}$ induce connected subgraphs.
\item Select the valid move that minimizes $\delta_{\max}$ (and equivalently $\Delta$).
\end{enumerate}

Both candidate moves preserve contiguity in both districts. Table~\ref{tab:c2_comparison} summarizes the district populations, balance metrics, and topological cut sizes before and after reassignment.

\begin{table}[htbp]
\centering
\caption{Comparison of district population balance and cut size before and after single boundary-block reassignment.}
\label{tab:c2_comparison}
\begin{tabular}{@{}lcccccc@{}}
\toprule
Partition Scenario & $P(G_1)$ & $P(G_2)$ & $\delta_{\max}$ (\%) & $\Delta$ (\%) & Cut Edges & Contiguous \\
\midrule
Baseline (Figure~1.22) & 1280 & 720 & 28.000 & 56.000 & 3 & Yes \\
Move block 2010 ($G_1 \to G_2$) & 1100 & 900 & 10.000 & 20.000 & 4 & Yes \\
\textbf{Move block 2004 ($G_1 \to G_2$)} & \textbf{1020} & \textbf{980} & \textbf{2.000} & \textbf{4.000} & \textbf{5} & \textbf{Yes} \\
\bottomrule
\end{tabular}
\end{table}

Reassigning block~2004 ($p_{2004} = 260$) from $G_1$ to $G_2$ yields the updated districts:
\begin{itemize}
\item $G_1' = \{2000, 2001, 2002, 2003, 2010, 2011, 2012, 2013, 2014, 2015\}$ with $P(G_1') = 1020$,
\item $G_2' = \{2004, 2005, 2006, 2007, 2008, 2009\}$ with $P(G_2') = 980$.
\end{itemize}

\textit{Interpretation:} As shown in Table~\ref{tab:c2_comparison}, moving block~2004 dramatically reduces the population difference from 560 to 40 residents, lowering $\delta_{\max}$ from $28.000\%$ to $2.000\%$ and bringing the overall deviation $\Delta$ down from $56.000\%$ to $4.000\%$---well within the $10\%$ legal equity standard. Both updated districts remain fully contiguous. However, this improvement increases the cut size from 3 to 5 severed edges ($(2000, 2004)$, $(2001, 2004)$, $(2011, 2004)$, $(2012, 2004)$, and $(2010, 2005)$), illustrating the fundamental trade-off in redistricting between population equity and geometric compactness.

\end{document}